\section{Arithmetic factorization and paired boxes}
\label{sec:arithmetic-boxes}

We convert a large prime-power divisor into an integral equation in a
fixed number field.  Balancing its factors at every embedding will put
their projective coordinates in compact sets.  The resulting equations
then pair small boxes in the two sets of coordinates, with only boundedly
many partners for each box.
This factorization-and-balancing strategy follows Reuss
\cite[Section~3]{Reuss2015PowerFree}, with the Gaussian factorization
in Heath-Brown \cite[Sections~2--4]{HeathBrown2012Squarefree} as an
earlier model. We use the standard ideal factorization, class-group
finiteness and full logarithmic unit lattice from
\cite[Theorems~3.7, 4.4 and~5.9]{milne2020ant}.

Throughout this section, $f\in\Z[T]$ is primitive and irreducible, has
degree $d$ and positive leading coefficient $c$, and $k=d-2\ge2$.
Fix a root $\theta$, put $K=\Q(\theta)$, and let $L$ be a splitting field
of $f$.  Write $\sigma_i:K\hookrightarrow L$ for the $d$ embeddings,
with $\theta_i=\sigma_i(\theta)$.  We consider
\[
 X<n\le2X,\qquad P\le p\le2P,\qquad p>X,\qquad p^k\mid f(n),
 \qquad P=X^\eta,
\]
and put $b=d-k\eta$.  All constants below may depend on $f$; none
depends on the individual $n,p$ or on $X$.

\begin{lemma}[Balanced factorization]\label{lem:balanced-factorization}
There are a positive integer $\mu$ and a fixed finite list of integral
ideals $\mathfrak c$ of $K$ with the following property.  For all
sufficiently large $X$, each pair $(n,p)$ above admits integral elements
$\alpha,\beta\in\mathcal O_K$ and one ideal from the list such that
\begin{equation}\label{eq:arithmetic-factorization}
 \alpha^k\beta=\mu(n-\theta),\qquad
 |N_{K/\Q}(\alpha)|=pN(\mathfrak c),
\end{equation}
and, for every embedding $\sigma_i$,
\begin{equation}\label{eq:arithmetic-balanced-sizes}
 |\sigma_i(\alpha)|\asymp_f P^{1/d},\qquad
 |\sigma_i(\beta)/\mu|\asymp_f X/P^{k/d}.
\end{equation}
We may select one representation for each $n$ under consideration.
For a fixed ideal case and fixed $\alpha$, there are at most $d$
selected values of $n$.
\end{lemma}

\begin{proof}
The element $c\theta$ is integral, with monic minimal polynomial
$F(T)=c^{d-1}f(T/c)\in\Z[T]$, and
\begin{equation}\label{eq:arithmetic-norm}
 N_{K/\Q}(cn-c\theta)=F(cn)=c^{d-1}f(n).
\end{equation}
Exclude the finitely many primes dividing $c$, the discriminant of $F$,
or $[\mathcal O_K:\Z[c\theta]]$; the factorization theorem for
this order is recalled in \cite[Theorem~3.41 and Remark~3.43]{milne2020ant}.
For every remaining prime $p$ dividing
$f(n)$, exactly one prime ideal $\mathfrak p$ above $p$ contains
$cn-c\theta$.  Indeed its residue factor must be the simple linear
factor $T-cn$ of $F$ modulo $p$.  Thus $N(\mathfrak p)=p$, and
\eqref{eq:arithmetic-norm} shows that its valuation in $(cn-c\theta)$
is exactly $v_p(f(n))$.  In particular,
$\mathfrak p^k\mid(cn-c\theta)$.  Every excluded prime is smaller than
$X$ once $X$ is sufficiently large.

Choose fixed integral representatives for the ideal classes of $K$.
One of them, denoted $\mathfrak c$, makes $\mathfrak p\mathfrak c$
principal.  Write
\[
 (\alpha)=\mathfrak p\mathfrak c.
\]
Then $\alpha$ is integral and its absolute norm is $pN(\mathfrak c)$.
We may multiply $\alpha$ by a unit to obtain the first estimate in
\eqref{eq:arithmetic-balanced-sizes}.  To see that this balances all
embeddings, use one logarithmic coordinate for each real embedding and
each complex pair, with respective weights $1$ and $2$.  Subtract
$d^{-1}\log(pN(\mathfrak c))$ from every coordinate.  The resulting
vector lies in the weighted sum-zero hyperplane.  Dirichlet's Unit
Theorem allows it to be moved into a bounded fundamental parallelepiped
of the unit lattice.  Exponentiating gives the asserted bounds at
every real and complex embedding.

Choose a positive integer $q$ such that
$q\mathfrak c^{-k}\subseteq\mathcal O_K$ for every ideal in the finite
list, and put $\mu=cq$.  Such a $q$ exists; for example one may take a
common multiple of the integers $N(\mathfrak c)^k$.  The ideal
$(cn-c\theta)\mathfrak p^{-k}$ is integral, so
\[
 \beta=\frac{\mu(n-\theta)}{\alpha^k}
\]
is integral.  Since $|n-\theta_i|\asymp_f X$, the second estimate in
\eqref{eq:arithmetic-balanced-sizes} follows from the first.

For each $n$ choose one admissible prime in the dyadic range, one ideal
case and one balanced generator; these choices determine $\beta$.
For a fixed ideal case, $\alpha$ determines
$p=|N_{K/\Q}(\alpha)|/N(\mathfrak c)$.  Outside the excluded primes,
$f$ has at most $d$ roots modulo $p^k$.  Since $p^k>X$, each such
residue class meets $(X,2X]$ in at most one integer.  This proves the
last assertion.
\end{proof}

Fix an integral basis $\omega_0,\ldots,\omega_{d-1}$ of $K$ and write
\[
 \alpha=\sum_{j=0}^{d-1}x_j\omega_j,\qquad
 \beta=\sum_{j=0}^{d-1}y_j\omega_j,\qquad x,y\in\Z^d.
\]
Define the conjugate linear forms
\[
 a_i(x)=\sum_j\sigma_i(\omega_j)x_j,\qquad
 b_i(y)=\mu^{-1}\sum_j\sigma_i(\omega_j)y_j.
\]
The selected points satisfy
\begin{equation}\label{eq:arithmetic-conjugate-equations}
 a_i(x)^k b_i(y)=n-\theta_i,\qquad
 |a_i(x)|\asymp_f X^{\eta/d},\qquad
 |b_i(y)|\asymp_f X^{b/d}.
\end{equation}
Both embedding matrices are invertible.  There is at most one selected
point for each $n$, and a fixed $x$ costs $O_f(1)$ selected points,
including all ideal cases.  These two properties survive any invertible
integral changes in $x$ and $y$.

\subsection{Uniform analytic coordinates}

Split into the finitely many choices of a largest absolute coordinate
of $x$ and of $y$, breaking ties in a fixed order.  Within each chart
reindex them as $x_0,y_0$.  Invertibility of the embedding matrices and
\eqref{eq:arithmetic-balanced-sizes} give
\[
 |x_0|\asymp_f P^{1/d},\qquad
 |y_0|\asymp_f X/P^{k/d}.
\]
Put
\[
 s_j=x_j/x_0,\qquad t_j=y_j/y_0\quad(1\le j<d),
 \qquad z=1/n.
\]
Thus $s,t\in[-1,1]^{d-1}$.  Every normalized conjugate form
$a_i(1,s)$ and $b_i(1,t)$ has modulus bounded above and below by
positive constants.  The chart vectors $s,t$ here are distinct from the
scalar parameter used later to choose bidegrees.

\begin{lemma}[Paired boxes]\label{lem:paired-boxes}
Fix $0<w<1$ and an integer $M\asymp X^w$.  Partition each of the
$s$- and $t$-cubes into grid boxes of side $M^{-1}$.  For the selected
points in any fixed arithmetic and coordinate chart, the following
hold uniformly.
\begin{enumerate}
\item The vector $t$ is locally an analytic function of $(s,z)$.
Conversely, $s$ is locally given by boundedly many analytic branches in
$(t,z)$.
\item The branches needed at a counted point extend to a neighborhood
of fixed radius containing the corresponding grid base with last
coordinate $z=0$.  Derivatives of every fixed order are uniformly
bounded there.
\item Every relevant box of either type has $O_f(1)$ partner boxes of
the other type.  In particular, there are $O_f(M^{d-1})$ relevant
pairs of boxes.
\end{enumerate}
When $m=(d-1)w$, the last bound is $O_f(X^m)$.  No primality assumption
on $M$ is needed for this lemma.
\end{lemma}

\begin{proof}
Let $\mathcal A$ and $\mathcal B$ be the invertible complex linear maps
from the chart coordinates to their $a$- and $b$-embedding vectors.
At a counted point put
\[
 \kappa=\frac{x_0^k y_0}{n}.
\]
The size bounds show $|\kappa|\asymp_f1$.  The normalized equations
\eqref{eq:arithmetic-conjugate-equations} give
\begin{equation}\label{eq:arithmetic-chart-forward}
 \left(\frac{1-\theta_i z}{a_i(1,s)^k}\right)_{i=1}^d
 =\kappa\,\mathcal B(1,t).
\end{equation}
Applying $\mathcal B^{-1}$ and dividing by its first coordinate gives
the forward formula for $t$.  Its normalizing denominator equals
$\kappa$ at the counted point, and is therefore bounded away from zero.

In the reverse direction, take local $k$-th roots of
\[
 \frac{1-\theta_i z}{b_i(1,t)}.
\]
The inputs lie in a fixed bounded annulus separated from zero.  At a
counted point, compatible roots have the form
$\gamma a_i(1,s)$ with $\gamma^k=\kappa$.  Applying
$\mathcal A^{-1}$ gives $\gamma(1,s)$; its first coordinate has modulus
$\asymp_f1$.  Thus division by that coordinate is again allowed.  There
are at most $k^d$ root choices, so only boundedly many branches are
needed.  Complex embeddings cause no additional restriction: we only
retain the branches that describe the counted real coordinate vectors.

These formulas also give uniform neighborhoods.  The normalized forms,
root inputs and the normalizing denominators have fixed upper bounds
and positive lower bounds at every counted point.  On a sufficiently
small complex polydisc, the same lower bounds hold with a smaller
constant.  Each root has a holomorphic branch on that polydisc, and
the rational operations remain holomorphic there.  Their formulas, or
Cauchy's estimates on a slightly smaller polydisc, bound every fixed
order of derivative uniformly.  This construction is local; it does
not require one choice of a root on the whole annulus.

The distance from a counted point to its grid base with $z=0$ is
$O(M^{-1}+X^{-1})$, so that base lies in the same fixed neighborhood
for large $X$.  The allowed $z$-interval has length $O(X^{-1})$.
Since $w<1$, this is $O(M^{-1})$.  The derivative bound therefore
maps one relevant $s$-box and the whole $z$-interval into a set of
$t$-diameter $O_f(M^{-1})$, meeting $O_f(1)$ grid boxes.  For a fixed
$t$-box, continue its relevant root choices to the common grid base.
There are at most $k^d$ such choices there.  Every branch realized by
a counted point has a uniformly nonzero normalizing denominator
throughout this small box and the allowed $z$-interval.  Each therefore
has an $s$-image of diameter $O_f(M^{-1})$, proving the reverse
partner bound.
There are $O(M^{d-1})$ boxes in either fixed cube, proving the total
bound.
\end{proof}

We have reduced the selected inputs to finitely many arithmetic cases
and paired projective boxes.  The equations
\eqref{eq:arithmetic-conjugate-equations} retain the integer input $n$,
while the analytic descriptions provide the small variables for the
determinant cuts in the next section.
